\documentclass[10pt,a4paper]{amsart}
\usepackage[margin=19mm]{geometry}
\usepackage{amsmath,amssymb,mathtools}
\usepackage[scr=rsfs,cal=euler]{mathalfa}
\usepackage{xfrac}
\usepackage[unicode,hidelinks,bookmarks=false]{hyperref}
\newcommand{\ie}{i.e.\ }
\newcommand{\cf}{cf.\ }
\newcommand{\resp}{respectively\ }
\newcommand{\Subsection}{\subsection}
\providecommand{\nobreakdash}{\nobreak\hbox{-}}
\setlength{\emergencystretch}{2em}
\multlinegap0pt
\usepackage{bm}
\usepackage{bbm}
\usepackage{dsfont}
\usepackage{xspace}
\usepackage{cancel}
\usepackage{mathtools}
\usepackage{amsthm}
\makeatletter
\@ifundefined{theorem}{\newtheorem{theorem}{Theorem}}{}
\@ifundefined{lemma}{\newtheorem{lemma}[theorem]{Lemma}}{}
\@ifundefined{proposition}{\newtheorem{proposition}[theorem]{Proposition}}{}
\makeatother
\newcommand{\At}{\operatorname{At}}
\newcommand{\CC}{\mathbb C}
\newcommand{\End}{\operatorname{End}}
\newcommand{\Hom}{\operatorname{Hom}}
\newcommand{\cHom}{\mathcal{H}om}
\newcommand{\cL}{\mathcal L}
\newcommand{\cO}{\mathcal O}
\newcommand{\tr}{\operatorname{tr}}
\title[Genus 4 Supermoduli Space Is Not Projected]{Genus 4 Supermoduli Space Is Not Projected}
\author{Ron Donagi, Simone Noja}
\date{Source: arXiv:2608.24996v2 (CC BY 4.0)}
\begin{document}
\maketitle

\noindent\emph{Source and licence.} Genus 4 Supermoduli Space Is Not Projected. Version arXiv:2608.24996v2 (CC BY 4.0), distributed under the Creative Commons Attribution 4.0 International licence, as recorded in the arXiv licence metadata for this exact version and shown on the abstract page https://arxiv.org/abs/2608.24996v2. Full rights notice: licenses/arxiv-2608.24996-CC-BY-4.0.txt.\par\smallskip

\noindent\textit{Editorial scope.} Counted unit: the Proposition whose source label is \texttt{prop:explicit-nonzero} in the article (source0.tex lines 1259--1266, source label \texttt{prop:explicit-nonzero}), delivered together with the article's own complete original proof of it (source0.tex lines 1268--1468, one \texttt{proof} environment of 201 lines). No mathematics is written, repaired or re-derived: statement and proof are the article's own text line for line. Required context retained uncounted: eq:j (lines 982--986); eq:En-extension (lines 1146--1152); eq:en-class-target (lines 1155--1159); eq:delta3 (lines 1195--1199); lem:delta-injective (lines 1217--1219); eq:End-identification (lines 1239--1241); eq:A-matrix (lines 1250--1254). Every definition, display and label of those lines is retained unchanged. Omitted from this package and not redistributed: 1--981, 987--1145, 1153--1154, 1160--1194, 1200--1216, 1220--1238, 1242--1249, 1255--1258, 1267--1267, 1469--1549. Nothing inside a retained excerpt is omitted or altered: every retained body is delivered exactly as the article's lines are written, so no omission receipt of any kind accompanies this package. Citations: the retained excerpts carry no citation command at all, so no bibliography entry of the article is transcribed and no reference is redistributed. Reference adaptation: none was needed and none was made; every reference inside the delivered unit resolves inside it, so no unresolved reference or citation is produced. Printed slips kept literal: no printed word, symbol, hypothesis or number of the counted statement or of its proof was altered. Layout and class: the article's own class is \texttt{amsart} with packages including inputenc, fontenc, babel, geometry, microtype, lmodern, mathpazo, euler, amsmath, amssymb, amsthm, mathtools, mathrsfs, bm; the wrapper is the lane's pinned \texttt{amsart} editorial wrapper at 10pt on A4, which restates the theorem-like environments the retained text needs and adds the article's own macro definitions that the retained text needs (\texttt{\textbackslash At}, \texttt{\textbackslash CC}, \texttt{\textbackslash End}, \texttt{\textbackslash Hom}, \texttt{\textbackslash cHom}, \texttt{\textbackslash cL}, \texttt{\textbackslash cO}, \texttt{\textbackslash tr}), with extra wrapper packages (bm, bbm, dsfont, xspace, cancel, mathtools). The delivered numbering is the unit's own. Assets: no retained span contains a figure environment, a graphics-inclusion command, a PDF-page inclusion, a table or a TikZ picture; no image file is copied into this package, the package declares no asset dependency and no image file is redistributed with it. Licence: version arXiv:2608.24996v2 of this article is distributed under the Creative Commons Attribution 4.0 International licence, as recorded in the arXiv licence metadata for this exact version and shown on the abstract page https://arxiv.org/abs/2608.24996v2. A full rights notice is delivered at licenses/arxiv-2608.24996-CC-BY-4.0.txt. Characters: every retained excerpt is pure ASCII, so the delivered engine, XeLaTeX with its default Latin Modern text font, reports no missing character.
\begin{equation}\label{eq:j}
 j:H^0(C,K_C)\otimes\cO_C\longrightarrow\cHom(E_2,E_3),
 \qquad
 s\longmapsto m_s,
\end{equation}

\begin{equation}\label{eq:En-extension}
0\longrightarrow V_n\otimes\cO_C
\longrightarrow E_n
\longrightarrow Q_n\longrightarrow0,
\qquad
Q_n=K_C^{n-1}\oplus\sigma^*K_C^{n-1}.
\end{equation}

\begin{equation}\label{eq:en-class-target}
 e_n\in
 \operatorname{Ext}^1_C(Q_n,V_n\otimes\cO_C)
 =H^1\!\left(C,\cHom(Q_n,V_n\otimes\cO_C)\right).
\end{equation}

\begin{equation}\label{eq:delta3}
 \delta_3:V_2\oplus V_2\longrightarrow\Hom(V_1,V_3),
 \qquad
 (q,q')\longmapsto N_q+DN_{q'}.
\end{equation}

\begin{lemma}\label{lem:delta-injective}
The map $\delta_3$ is injective.
\end{lemma}

\begin{equation}\label{eq:End-identification}
H^1(C,\cO_C)\otimes H^0(C,K_C)\simeq\End(V_1).
\end{equation}

\begin{equation}\label{eq:A-matrix}
 A:=\At_{\mathrm{mix}}(\cL)
 =I-\rho_1^{-1}
 =\operatorname{diag}(-\zeta,-\zeta^2).
\end{equation}

\begin{proposition}\label{prop:explicit-nonzero}
For the multiplication morphism \eqref{eq:j},
$$
 H^1(j)(A)\neq0
 \quad\text{in}\quad
 H^1\!\left(C,\cHom(E_2,E_3)\right).
$$
\end{proposition}

\begin{proof}
Choose a sufficiently fine \v Cech cover
$\mathfrak U=\{U_a\}$ of $C$ on which both extensions
\eqref{eq:En-extension} split.  Relative to local splittings
$$
 E_n|_{U_a}\simeq
 (V_n\otimes\cO_{U_a})\oplus Q_n|_{U_a},
$$
the comparison maps between these splittings are automorphisms of
the restrictions of the fixed bundles
$(V_n\otimes\cO_C)\oplus Q_n$ to $U_{ab}$.
Writing $\chi_{n,a}:(V_n\otimes\cO_{U_a})\oplus Q_n|_{U_a}
\xrightarrow{\sim}E_n|_{U_a}$ for these identifications, the comparison
$T_{n,ab}=\chi_{n,a}^{-1}\chi_{n,b}$ has matrix
\begin{equation}\label{eq:Tnab}
 T_{n,ab}=
 \begin{pmatrix}I_{V_n}&e_{n,ab}\\0&I_{Q_n}\end{pmatrix}.
\end{equation}
Here
$$
 e_{n,ab}\in
 \Gamma\!\left(U_{ab},
 \cHom(Q_n,V_n\otimes\cO_C)\right),
$$
and the \v Cech cocycle $(e_{n,ab})$ represents the extension class
$e_n$ in \eqref{eq:en-class-target}.

Choose a \v Cech cocycle
$\alpha=(\alpha_{ab})\in
Z^1(\mathfrak U,V_1\otimes\cO_C)$ representing the class $A$ under
\eqref{eq:End-identification}.  Applying the sheaf morphism $j$ gives
$$
 Z_A:=j(\alpha)\in
 Z^1\!\left(\mathfrak U,\cHom(E_2,E_3)\right).
$$
In the splittings \eqref{eq:Tnab}, each local homomorphism
$Z_{A,ab}:E_2|_{U_{ab}}\to E_3|_{U_{ab}}$ has the form
\begin{equation}\label{eq:ZA-blocks}
 Z_{A,ab}=
 \begin{pmatrix}
  Z_{A,ab}^{\mathrm{ul}}&Z_{A,ab}^{\mathrm{ur}}\\
  0&Z_{A,ab}^{\mathrm{lr}}
 \end{pmatrix}.
\end{equation}
The lower-left block is zero because multiplication preserves the constant
subbundles $V_2\otimes\cO_C\subset E_2$ and
$V_3\otimes\cO_C\subset E_3$.  More explicitly,
$$
\begin{aligned}
 Z_{A,ab}^{\mathrm{ul}}
 &\in\Gamma\!\left(
 U_{ab},\cHom(V_2,V_3)\otimes\cO_C\right),\\
 Z_{A,ab}^{\mathrm{lr}}
 &\in\Gamma\!\left(U_{ab},\cHom(Q_2,Q_3)\right).
\end{aligned}
$$

We must determine whether the equation
\begin{equation}\label{eq:solve-coboundary}
 Z_A=\delta\Phi
\end{equation}
has a solution
$\Phi\in C^0(\mathfrak U,\cHom(E_2,E_3))$; this is precisely the equation
whose solvability would say that $H^1(j)(A)=0$.  Write a putative local
zero-cochain as
\begin{equation}\label{eq:Phi-blocks}
 \Phi_a=
 \begin{pmatrix}
  \phi_a&\beta_a\\
  \gamma_a&\psi_a
 \end{pmatrix}.
\end{equation}
The four entries in \eqref{eq:Phi-blocks} are, respectively, local sections
of
$$
\begin{gathered}
 \cHom(V_2,V_3)\otimes\cO_C,\qquad
 \cHom(Q_2,V_3\otimes\cO_C),\\
 \cHom(V_2\otimes\cO_C,Q_3),\qquad
 \cHom(Q_2,Q_3).
\end{gathered}
$$
We use the coboundary convention
$$
 (\delta\Phi)_{ab}
 =T_{3,ab}\Phi_bT_{2,ab}^{-1}-\Phi_a.
$$
Multiplication of these block matrices gives
\begin{align}
 (\delta\Phi)_{ab}^{\mathrm{ll}}
 &=\gamma_b-\gamma_a,\label{eq:delta-ll}\\
 (\delta\Phi)_{ab}^{\mathrm{ul}}
 &=\phi_b-\phi_a+e_{3,ab}\gamma_b,\label{eq:delta-ul}\\
 (\delta\Phi)_{ab}^{\mathrm{lr}}
 &=\psi_b-\psi_a-\gamma_be_{2,ab}.\label{eq:delta-lr}
\end{align}
Since the lower-left block of $Z_A$ vanishes,
\eqref{eq:solve-coboundary} and \eqref{eq:delta-ll} force the
$\gamma_a$ to glue to a global morphism
$$
 \gamma:V_2\otimes\cO_C\longrightarrow Q_3.
$$
Taking cohomology in the upper-left and lower-right blocks yields the two
necessary identities
\begin{align}
 [Z_A^{\mathrm{ul}}]
 &=[e_3\gamma]
 &&\text{in }H^1\!\left(
 C,\cHom(V_2,V_3)\otimes\cO_C\right),
 \label{eq:block-identity-ul}\\
 [Z_A^{\mathrm{lr}}]
 &=-[\gamma e_2]
 &&\text{in }H^1\!\left(C,\cHom(Q_2,Q_3)\right).
 \label{eq:block-identity-lr}
\end{align}

We first solve the upper-left identity.  Using
$H^0(C,Q_3)\simeq V_2\oplus V_2$, write the global morphism $\gamma$ as
\begin{equation}\label{eq:gamma-C1C2}
 \gamma=(C_1,C_2):V_2\longrightarrow V_2\oplus V_2,
 \qquad C_1,C_2\in\End(V_2).
\end{equation}
Here $(C_1,C_2)$ denotes the induced map on global sections;
the bundle morphism $\gamma$ is obtained from it by evaluation into $Q_3$.
Under Serre duality, evaluation of
$[Z_A^{\mathrm{ul}}]$ on $q\in V_2$ is the map
$N_qA\in\Hom(V_1,V_3)$.  Evaluation of $[e_3\gamma]$ is the connecting
map \eqref{eq:delta3} applied to $(C_1q,C_2q)$.  Therefore
\eqref{eq:block-identity-ul} is equivalent to
\begin{equation}\label{eq:upper-equation}
 N_qA=N_{C_1q}+DN_{C_2q}
 \qquad\text{for every }q\in V_2.
\end{equation}
Equivariance of multiplication gives
$\rho_3N_q=N_{\rho_2q}\rho_1$, and therefore
$$
 DN_{\rho_2q}=N_q\rho_1^{-1}.
$$
Since $A=I-\rho_1^{-1}$ by \eqref{eq:A-matrix}, the pair
\begin{equation}\label{eq:C1C2}
 C_1=I_3,
 \qquad
 C_2=-\rho_2
\end{equation}
solves \eqref{eq:upper-equation}.  It is the unique solution because
$\delta_3$ is injective by Lemma~\ref{lem:delta-injective}.

It remains to test the lower-right identity
\eqref{eq:block-identity-lr}.  Using the differential of $\sigma$ to
identify the second summands, we have
$Q_2\simeq K_C\oplus K_C$ and
$Q_3\simeq K_C^2\oplus K_C^2$.  Hence
\begin{equation}\label{eq:lower-target}
 H^1\!\left(C,\cHom(Q_2,Q_3)\right)
 \simeq M_2\!\left(H^1(C,K_C)\right)
 \simeq M_2(\CC).
\end{equation}
The row and column order in this matrix space is the order
$(\Delta,\Gamma_\sigma)$ of the two quotient summands.

Use the Serre trace $H^1(C,K_C)\xrightarrow{\sim}\CC$ in
\eqref{eq:lower-target}.  Contraction of dual Serre bases with a canonical
extension tensor is then the ordinary matrix trace.  On the two
quotient summands, multiplication by $s\in V_1$ is
$\operatorname{diag}(s,\rho_1s)$.  Therefore the lower-right class of the
cocycle $Z_A$ is
\begin{equation}\label{eq:ZA-lr}
 [Z_A^{\mathrm{lr}}]
 =\operatorname{diag}\bigl(\tr A,\tr(A\rho_1)\bigr)
 =\operatorname{diag}(1,-1)
 \quad\text{in }M_2(\CC).
\end{equation}
On the other hand, write $R_1=I$ and $R_2=\rho_2^{-1}$ for the
two tensors of $e_2$.  For dual Serre bases $(q_k)$ and $(q_k^*)$, the
$(i,j)$ entry of $[\gamma e_2]$ is
$$
 \sum_k q_k^*(C_iR_jq_k)=\tr(C_iR_j).
$$
Thus composing \eqref{eq:C1C2} with the extension tensor
$e_2=(I,\rho_2^{-1})$ gives
\begin{equation}\label{eq:gamma-e2}
 [\gamma e_2]=
 \begin{pmatrix}
  \tr I_3&\tr\rho_2^{-1}\\
  -\tr\rho_2&-\tr I_3
 \end{pmatrix}
 =\operatorname{diag}(3,-3),
\end{equation}
because $1+\zeta+\zeta^2=0$.  Thus the residual lower-right class is
\begin{equation}\label{eq:final-matrix}
 [Z_A^{\mathrm{lr}}+\gamma e_2]
 =\begin{pmatrix}4&0\\0&-4\end{pmatrix}\neq0
 \quad\text{in }H^1\!\left(C,\cHom(Q_2,Q_3)\right).
\end{equation}
This contradicts the necessary identity
\eqref{eq:block-identity-lr}.  The lower-left, upper-left, and lower-right
blocks already give incompatible necessary conditions, so no equation
from the upper-right block is needed.  Consequently
\eqref{eq:solve-coboundary} has no solution, and
$H^1(j)(A)\neq0$.
\end{proof}

\end{document}
